\chapter*{译者说明}
\addcontentsline{toc}{chapter}{译者说明}
\markboth{译者说明}{译者说明}

\begin{flushleft}

本书依据 \textit{\booktitleoriginal}（\bookauthororiginal\ 著）译出。译文力求贴近原文的论证层次与语气，同时按中文的阅读习惯调整句法；凡涉及作者核心概念、历史背景或双关之处，以脚注或“译者注”说明。

\vspace{0.6em}
\textbf{关于术语。}全书人名、地名、机构名与学术术语原则上从通行译法；同一概念只用一个译名。首次出现时以“中文（原文）”标出，之后径用中文。外文斜体词在译文中保留原貌的，多为作者作为分析对象反复使用的概念。

\vspace{0.6em}
\textbf{关于注释。}原书脚注全部译出，编号与原文对应；译者所加注释以“译者注”标明，不与原注混同。原文中的页码、年份、档案编号一律照录。

\vspace{0.6em}
\textbf{关于引文。}长引文独立排印，短引文随文加引号。引文内的省略号、方括号补充均按原文处理，不擅自补足或删节。

\vspace{0.6em}
翻译中的取舍难免有疏漏，欢迎读者指正。勘误将在后续版本中修订。

\vspace{1.2em}
\hfill 译者

\end{flushleft}
